\begin{deluxetable}{lccc}
\tablecaption{Per-Line Measurements\label{table:linesummary}}
\tablewidth{0pt}
\tablehead{
  \colhead{Symbol} &
  \colhead{Column\tablenotemark{a}} &
  \colhead{Units} &
  \colhead{Inverse variance?}
}
\startdata
$A$                       & \texttt{MODELAMP}    & \fluxdensity & \nodata \\
$A_\mathrm{obs}$          & \texttt{AMP}         & \fluxdensity & Y \\
$\Phi$                    & \texttt{FLUX}        & \flux        & Y \\
$\Phi_\mathrm{box}$       & \texttt{BOXFLUX}     & \flux        & Y \\
$\Delta v$                & \texttt{VSHIFT}      & \kms         & Y \\
$\sigma$                  & \texttt{SIGMA}       & \kms         & Y \\
$F_c$                     & \texttt{CONT}        & \fluxdensity & Y \\
EW                        & \texttt{EW}          & \AA          & Y \\
$\Phi_\mathrm{lim}$       & \texttt{FLUX\_LIMIT} & \flux        & \nodata \\
$\chi^2_\mathrm{line}$    & \texttt{CHI2}        & \nodata      & \nodata \\
$n_\mathrm{pix}^\mathrm{line}$ & \texttt{NPIX}   & \nodata      & \nodata \\
\enddata
\tablenotetext{a}{Each column is prefixed by the line name (e.g.,
\texttt{OIII\_5007\_FLUX}); Table~\ref{table:emlines} lists the full
set of modeled lines.}
\tablecomments{Columns marked ``Y'' have a companion \texttt{\_IVAR}
column storing the inverse variance, estimated via the same
per-object Monte Carlo procedure described in \S\ref{ssec:uncertainties}.
$A$, $\Delta v$, and $\sigma$ are defined in \S\ref{sssec:emline_model};
$A_\mathrm{obs}$, $\Phi$, $\Phi_\mathrm{box}$, $F_c$, EW, $\Phi_\mathrm{lim}$,
$\chi^2_\mathrm{line}$, and $n_\mathrm{pix}^\mathrm{line}$ are defined in
\S\ref{sssec:lines}.}
\end{deluxetable}
